% Problema, público interno o externo involucrado y KPI(s) de impacto esperado

\section{Resumen Ejecutivo}

El presente proyecto aborda la detección de fraude en transacciones financieras mediante técnicas de aprendizaje automático supervisado. A partir del dataset sintético \textit{BankSim}, se desarrollarán y compararán modelos de clasificación capaces de asignar un nivel de riesgo a las operaciones, con el objetivo de facilitar la identificación de clientes potencialmente comprometidos.

El proyecto se orienta a un escenario de uso por parte de equipos de prevención de fraude de entidades financieras, cuya capacidad de revisión de alertas es limitada. Por este motivo, se utilizará como principal KPI \textit{Customer Precision@100}, que medirá la proporción de clientes realmente comprometidos entre los 100 clientes con mayor riesgo identificados diariamente. El desempeño se comparará con una regresión logística utilizada como \textit{baseline}, complementando la evaluación mediante Average Precision y ROC-AUC.